% This file was converted to LaTeX by Writer2LaTeX ver. 1.9.9
% see http://writer2latex.sourceforge.net for more info
\documentclass{article}
\usepackage{calc,amsmath,amssymb,amsfonts}
\usepackage[T1]{fontenc}
\usepackage[italian]{babel}
\usepackage[style=numeric,backend=biber]{biblatex}
\author{utente}
\date{2026-08-20}
\begin{document}
Quarant'anni fa, davanti a un Commodore 64, quattro ragazzi immaginarono di creare un videogioco tutto loro. Avevano
entusiasmo, curiosità e una macchina che sembrava capace di qualsiasi cosa. Quel gioco, però, non venne mai realizzato.

Oggi Vincenzo Camporeale torna davanti allo stesso computer per riprendere quel progetto rimasto incompiuto e provare
finalmente a portarlo a termine.

Questo libro racconta il viaggio.

Non è soltanto un ritorno nostalgico agli anni Ottanta e non è un tradizionale manuale di programmazione. È il diario
ragionato della costruzione di un vero videogioco per Commodore 64, sviluppato passo dopo passo, lasciando che siano i
problemi incontrati lungo il percorso a introdurre le tecniche necessarie per risolverli.

All'origine di quel percorso c'è anche un insegnante, il professor Rossi Brunori, che nei primi anni Ottanta portava
personalmente il proprio Commodore 64 in una normale aula scolastica e, davanti a un gruppo di ragazzi raccolti intorno
a quell'unica macchina, non si limitava a insegnare il BASIC: insegnava soprattutto a ragionare. Molte delle domande
nate durante quelle lezioni accompagneranno l'autore ben oltre gli anni della scuola, fino a trasformare quella prima
curiosità in una professione. E la storia tra l'allievo e il suo professore, cominciata davanti a un cursore
lampeggiante, attraverserà idealmente tutto il libro fino a trovare, molti anni dopo, un epilogo capace di chiudere un
cerchio rimasto aperto per oltre quarant'anni.

Nel frattempo il videogioco prende forma.

Si parte dal BASIC, dagli sprite e dai primi movimenti sullo schermo. Poi il progetto cresce e, con lui, cresce la
necessità di comprendere sempre più profondamente la macchina: joystick, collisioni, grafica, VIC-II, suono, SID,
raster, interrupt, organizzazione della memoria e, quando il BASIC non basta più, Assembly.

Ogni nuova funzionalità diventa così l'occasione per entrare un po' più a fondo nel funzionamento del Commodore 64. Non
vengono presentate soluzioni già confezionate: prima nasce un problema, poi si formulano ipotesi, si sperimenta, si
sbaglia, si osservano i limiti della soluzione adottata e soltanto allora si cerca una strada migliore. Perché
programmare non significa soltanto conoscere istruzioni e linguaggi, ma imparare a trasformare un'idea in qualcosa che
una macchina possa realmente eseguire.

A rendere questo viaggio ancora più singolare è il confronto tra due epoche tecnologiche lontanissime. Accanto al
Commodore 64, una delle macchine che negli anni Ottanta rappresentavano il futuro, oggi compare uno strumento che
allora sarebbe sembrato fantascienza: l'intelligenza artificiale. ChatGPT accompagna la progettazione e il ragionamento
come interlocutore contemporaneo, senza sostituire lo studio, le decisioni, gli esperimenti e gli errori che restano al
centro dello sviluppo.

È un libro per chi ha conosciuto il Commodore 64 e vuole tornare a guardarlo con occhi diversi; per chi ci ha giocato
per anni senza aver mai scoperto davvero cosa accadesse dietro lo schermo; per chi ha scritto i propri primi programmi
davanti a quella schermata azzurra; ma anche per chi, pur non avendo vissuto quell'epoca, vuole capire perché una
macchina con appena 64 kilobyte di memoria sia riuscita a insegnare a un'intera generazione un modo particolare di
pensare l'informatica.

Alla fine, il viaggio che separa un semplice PRINT da un videogioco completo nasce dalla stessa curiosità che, molti
anni fa, spinse un bambino a osservare una pallina rimbalzare sullo schermo di Pong e a porsi una domanda tanto
semplice quanto destinata a cambiargli la vita:

Come fa?


\bigskip
\end{document}
